\section{Exact equivariance of the arithmetic theta family}
\label{te-section}

An eighth-power identity does not by itself give an identity of
theta functions with the prescribed scalar.  The arithmetic
hypotheses supply the additional information needed here.  We
prove the exact identities, including the determinant characters.
The argument uses unramifiedness of the CM extension at every
finite place and the even dimensions in the construction.

\subsection{The trace lattices and their semicharacter}

Let $k=\mathbb Q(\sqrt{-D})$, $\sigma=\sqrt{-D}=i\sqrt D$,
and $\tau=(1+\sigma)/2$, with
$\mathcal O_k=\mathbb Z[\tau]$ and $D\equiv7\pmod8$.
Let $E/F$ be a CM quadratic extension unramified at every
finite place, with $k\subset E$ and $F$ its maximal real
subfield.  Write $\mathcal O=\mathcal O_E$ and let $\Sigma$
be the embeddings of $E$ into $\mathbb C$ extending the
specified embedding of $k$; put $d=|\Sigma|$.
Choose a fractional $\mathcal O$-ideal $\mathfrak a$ with
\[
          \mathfrak a\overline{\mathfrak a}
                         =\mathfrak D_{E/k}^{-1}.
 \tag{TE1}
\]
This is the ideal denoted by the inverse of the chosen half
different in the construction.  Its existence is part of the
arithmetic datum; the assertion below retains the exact equality
(TE1).  Let $I$ be the elliptic ideal of norm $a$, and let
$\lambda\in\mathbb C^\times$ be its common complex scaling.

For unimodular Hermitian $\mathcal O$-lattices
$L\subset E^n$, $M\subset E^{n'}$, define
\[
 \begin{split}
 W_{n,n'}&=\prod_{\eta\in\Sigma}
                      \operatorname{Mat}_{n'\times n}(\mathbb C),\\
 \Lambda(L,M)&=\{(\lambda\eta(X))_{\eta\in\Sigma}:
             X\in I\mathfrak a\operatorname{Hom}_{\mathcal O}(L,M)\},\\
 H(X,Y)&=\frac{2}{a|\lambda|^2\sqrt D}
                  \sum_{\eta\in\Sigma}\operatorname{tr}(X_\eta Y_\eta^*),
       \qquad E_H=\operatorname{Im}H.
 \end{split}
 \tag{TE2}
\]
The symbol $E_H$ denotes the alternating form; the unadorned
$E$ continues to denote the number field.  The standard
Hermitian forms on the two vector spaces are used in (TE2).
Put $L_0=\mathcal O^n$ and $M_0=\mathcal O^{n'}$.

\begin{lemma}\label{te-trace-lattice}
Every lattice in (TE2) is self-dual for $E_H$.
All these lattices are commensurable and stable under
$\mathcal O$.  On each of them
\[
       Q(x)=\frac{\sqrt D}{2}H(x,x)\in\mathbb Z,
                    \qquad \alpha(x)=e^{\pi iQ(x)}\in\{1,-1\}
 \tag{TE3}
\]
is the same sign semicharacter.
\end{lemma}
\begin{proof}
The Hermitian trace dual of
$\operatorname{Hom}_{\mathcal O}(L,M)$ for the $E$-valued
form $\operatorname{tr}(XY^*)$ is itself.  To check this,
localize at a prime, use a basis of $L$ and a basis of $M$,
and use their Hermitian dual bases at the conjugate prime.
The matrix units and their duals give precisely the asserted
dual module.  This argument does not assume that the global
projective modules have bases.

For the $k$-valued trace form, taking a dual inserts
$\mathfrak D_{E/k}^{-1}$.  Thus the dual of
$I\mathfrak a\operatorname{Hom}(L,M)$ for that form divided
by $a$ is
\[
 a\overline I^{-1}\overline{\mathfrak a}^{-1}
       \mathfrak D_{E/k}^{-1}\operatorname{Hom}(L,M)
             =I\mathfrak a\operatorname{Hom}(L,M).
\]
Here $I\overline I=(a)$ and (TE1) were used.
The inverse different of $k/\mathbb Q$ is
$\sigma^{-1}\mathcal O_k$.  Since
\[
 \operatorname{Tr}_{k/\mathbb Q}(t/\sigma)
                              =2\operatorname{Im}(t)/\sqrt D,
\]
passing to the alternating form in (TE2) gives exactly
the same dual lattice.  The factors $\lambda$ cancel in
the numerator and denominator.  This proves self-duality.
Full lattices in the same rational vector space are
commensurable, and stability under $\mathcal O$ is immediate.
Finally $Q(x)=E_H(\tau x,x)$, so integrality and the
semicharacter law follow from (TC42).
\end{proof}

Use the normalized averaging operators $I_{\Lambda,\Lambda'}$
of (TC10), with positive square roots of their indices.
No change in that normalization will be made below.

\subsection{A common local hyperbolic decomposition}

\begin{lemma}\label{te-local-pair}
Let $A$ be a complete discrete valuation ring of
characteristic zero with finite residue field, and let $B$
be its unramified quadratic extension, with involution
$x\mapsto\bar x$.  Let $V$ be an even-dimensional
nondegenerate Hermitian space over $\operatorname{Frac}(B)$.
If $M,N$ are unimodular $B$-lattices in $V$, there is a
decomposition into hyperbolic planes with vectors $v_j,u_j$
and integers $a_j\geq0$ such that
\[
 \begin{split}
 M&=\bigoplus_j(Bv_j\oplus Bu_j),\\
 N&=\bigoplus_j(B\varpi^{-a_j}v_j\oplus B\varpi^{a_j}u_j),\\
 h(v_i,v_j)&=h(u_i,u_j)=0,\qquad h(v_i,u_j)=\delta_{ij}.
 \end{split}
 \tag{TE4}
\]
In particular, an idempotent $e$ preserves both lattices
and satisfies $e^*=1-e$.
For a split quadratic algebra $B=A\times A$, an idempotent
with the same properties exists in every dimension.
\end{lemma}
\begin{proof}
We use the convention that $h$ is conjugate-linear in its
first variable.  Trace and norm on units are surjective by
Lemma~\ref{acl-local-norm}, including in residue
characteristic two.

First a single unimodular even-dimensional lattice is a
sum of hyperbolic planes.  Proposition~\ref{acl-local-basis}
gives an orthonormal basis.  In any pair of its coordinates
choose $t\in B^\times$ with $t\bar t=-1$; the vector
$(1,t)$ is primitive and isotropic.  Unimodularity gives
a vector $z$ paired with it to one.  Replacing $z$ by
$z-c(1,t)$, where $c+\bar c=h(z,z)$, makes $z$ isotropic
without changing this pairing.  The resulting hyperbolic
plane splits integrally, and induction proves the assertion.

For two lattices, let $a\geq0$ be the least integer with
$\varpi^aN\subset M$.  Duality also gives
$\varpi^aM\subset N$.  If $a=0$, the lattices are equal
and the preceding argument applies.  If $a>0$, choose
$w\in\varpi^aN$ primitive in $M$; otherwise $a$ was not
minimal.  Its norm lies in $\varpi^{2a}A$.
Choose $z\in M$ with $h(w,z)=1$ and choose $c\in B$
whose trace is a unit.  The equation
\[
 h(w+\varpi^{2a}ctz,w+\varpi^{2a}ctz)=0
                    \qquad(t\in A)
 \tag{TE5}
\]
divided by $\varpi^{2a}$ has the form
$b+t\operatorname{Tr}(c)+\varpi^{2a}N(c)h(z,z)t^2=0$.
It has a root modulo $\varpi$, and its derivative is a
unit there.  Successively correcting the root modulo
$\varpi^r$ and taking its limit gives a root in $A$.
This also supplies the required elementary form of
Hensel's argument.

The resulting vector $v$ is primitive in $M$, is isotropic,
and still belongs to $\varpi^aN$, because
$\varpi^{2a}M\subset\varpi^aN$.
Choose $u\in M$ with $h(v,u)=1$ and make it isotropic
by subtracting $cv$ with $c+\bar c=h(u,u)$.
The vectors $v,u$ split a unimodular hyperbolic plane
off $M$.  The vectors $\varpi^{-a}v,\varpi^a u$
belong to $N$, are isotropic, and have pairing one;
they split the same plane off $N$.
The two orthogonal complements are again unimodular
lattices of even rank.  Induction gives (TE4).
Projection onto the span of the $v_j$ has the stated
properties.  In the split case use multiplication by
$(1,0)$; its adjoint is multiplication by $(0,1)$.
\end{proof}

\begin{proposition}\label{te-projector}
Let $M,N$ be unimodular $\mathcal O$-lattices of even
rank in one Hermitian $E$-space.  For every positive
rational integer $b$ there is an $E$-linear map $e$
preserving both lattices such that each of
\[
 e^2-e,\quad e+e^*-1,\quad e^*e,\quad ee^*,\quad
                 (1-e)^*(1-e),\quad(1-e)(1-e)^*
 \tag{TE6}
\]
maps $M+N$ into $b(M\cap N)$.
\end{proposition}
\begin{proof}
Put $\mathscr R=\operatorname{End}_{\mathcal O}(M)
\cap\operatorname{End}_{\mathcal O}(N)$.  It is a full
finitely generated $\mathbb Z$-lattice in
$\operatorname{End}_E(V)$ and a ring.  It is stable under
adjoint, since $M$ and $N$ are self-dual.
Choose a positive integer $c$ with
$c(M+N)\subset b(M\cap N)$.
At each finite place of $F$ above a divisor of $c$,
Lemma~\ref{te-local-pair} supplies an exact idempotent
preserving both completed lattices and satisfying
$e_v^*=1-e_v$.  All polynomials in (TE6) vanish at $e_v$.

Completion commutes with this intersection of two full
lattices: after clearing denominators it is the kernel
of a homomorphism to a finite module, and completion is
exact on finitely generated modules over $\mathbb Z$.
Thus these local maps belong to the corresponding
completions of $\mathscr R$.
The Chinese remainder theorem on the finite free
$\mathbb Z$-module $\mathscr R$ gives
$e\in\mathscr R$ congruent to all of them modulo
$c\mathscr R$ at the primes dividing $c$.
Every expression in (TE6) then lies in $c\mathscr R$;
membership is checked at these primes, and elsewhere
$c$ is a unit.  Its image of $M+N$ lies in
$c(M+N)\subset b(M\cap N)$, as claimed.
\end{proof}

\subsection{The finite quadratic sum has positive sign}

For three commensurable self-dual lattices
$\Lambda_1,\Lambda_2,\Lambda_3$, recall the finite group
and its quadratic function from the proof of
Theorem~\ref{tc-triple}:
\[
 \begin{split}
 P&=\{(z,y)\in\Lambda_3\times\Lambda_2:z+y\in\Lambda_1\},\\
 T&=\{(t,u-t):t\in\Lambda_2\cap\Lambda_3,
                              \ u\in\Lambda_1\cap\Lambda_2\},\\
 \psi(z,y)&=\alpha(z)\alpha(y)\alpha(z+y)e^{-\pi iE_H(z,y)}.
 \end{split}
 \tag{TE7}
\]
The function $\psi$ is a well-defined homogeneous quadratic
function on $P/T$.  Its sum differs from the scalar
in $I_{23}I_{12}=\mu I_{13}$ by a positive real factor,
and $|\mu|=1$.

\begin{lemma}\label{te-positive-gauss}
Suppose an endomorphism $e$ of the ambient rational space
preserves the three lattices, commutes with $k$, and
each of $e^2-e$, $e^*e$, and $(1-e)^*(1-e)$ carries
$\Lambda_1+\Lambda_2+\Lambda_3$ into
$\Lambda_1\cap\Lambda_2\cap\Lambda_3$.
Then $\mu=1$.
\end{lemma}
\begin{proof}
The map $(z,y)\mapsto(ez,ey)$ induces an idempotent on
$P/T$: its square minus itself has both coordinates in
the common intersection, and every such pair belongs
to $T$.  Consequently
$P/T=A\oplus B$, where $A$ and $B$ are the images of
$e$ and $1-e$.

Put
\[
 \beta(x,y)=\sqrt D\operatorname{Re}H(x,y)-E_H(x,y)
                         =2E_H((\tau-1)x,y).
 \tag{TE8}
\]
If $x,y,x+y$ lie in the indicated three lattices,
integrality of their $Q$-values gives
$\psi(x,y)=e^{\pi i\beta(x,y)}$.
For $(z,y)\in P$,
\[
       \beta(ez,ey)=\beta(z,e^*ey)\in2\mathbb Z,
\]
because $z,e^*ey\in\Lambda_3$ and $\Lambda_3$ is
$\mathcal O_k$-stable.  Hence $\psi=1$ on $A$.
The same argument with $1-e$ proves $\psi=1$ on $B$.
Writing $b_\psi$ for the polar bicharacter gives
\[
       \sum_{x\in P/T}\psi(x)
                    =\sum_{a\in A}\sum_{b\in B}b_\psi(a,b)\geq0.
 \tag{TE9}
\]
Indeed the inner sum of any character of a finite group
is either its order or zero, as follows by translating
the sum by an element where the character is nontrivial.
The scalar $\mu$ is therefore nonnegative real.  Its
absolute value one implies $\mu=1$.
\end{proof}

\begin{proposition}\label{te-exact-paths}
Let $n,n'$ be even.  Normalized averaging satisfies both
triangle identities with one of $L,M$ fixed, and the
square identity
\[
 \begin{split}
 I_{\Lambda(L_1,M),\Lambda(L_1,M_1)}
 I_{\Lambda(L,M),\Lambda(L_1,M)}
 \;=\;
 I_{\Lambda(L,M_1),\Lambda(L_1,M_1)}
 I_{\Lambda(L,M),\Lambda(L,M_1)}
 \end{split}
 \tag{TE10}
\]
on the theta line of $\Lambda(L,M)$.
\end{proposition}
\begin{proof}
For the square, choose $b$ with
$b(L+L_1)\subset L\cap L_1$, and apply
Proposition~\ref{te-projector} to $M,M_1$.
The induced map on matrices is $X\mapsto eX$.
It preserves each of the four lattices.
Every error in (TE6) sends each of them into their
common intersection.  In fact, if $A:L_i\to M_j$
and $v\in L_k$, then $bv\in L_i$ and
\[
       R A(v)=b^{-1}R A(bv)\in M\cap M_1
\]
for each error $R$, since $R(M+M_1)\subset b(M\cap M_1)$.
Multiplication by $I\mathfrak a$ preserves this inclusion.
Lemma~\ref{te-positive-gauss} applies to either of the
two triples from the lower left to the upper right
corner of the square.  Both composites are the same
direct averaging operator, proving (TE10).

For three varying domain lattices and a fixed target
$M$, choose $b$ carrying their sum into their
intersection, and apply the same argument to the pair
$M,M$.  This proves the left triangle identity.
For three varying target lattices and fixed $L$, use
Proposition~\ref{te-projector} for $L,L$, with $b$
carrying the sum of those target lattices into their
intersection.  Use the induced map $X\mapsto Xe^*$.
Its adjoint is $X\mapsto Xe$; the errors are precisely
the corresponding right multiplications by the
expressions in (TE6) or their adjoints.  They have the
required images by the same calculation.  This proves
the right triangle identity.  Zero-dimensional cases
have identity averaging and satisfy all the formulas.
\end{proof}

\subsection{The determinant in the unitary action}

We next remove the possible sign from the action on an
even number of copies of a theta line.

\begin{lemma}\label{te-negative-norm}
For every positive rational integer $m$ there is
$\ell\in\mathcal O$ with
\[
                         \ell\bar\ell\equiv-1\pmod{m\mathcal O}.
 \tag{TE11}
\]
\end{lemma}
\begin{proof}
At an inert place of $F$ dividing $m$, choose a unit
of norm $-1$ by Lemma~\ref{acl-local-norm}; at a split
place use $(1,-1)$.  Approximate all these units modulo
$m$ in the product of the completed integer rings.
The Chinese remainder theorem for $\mathcal O$ gives
an element $\ell$ with these residues.  Multiplication
and conjugation preserve the congruences, proving (TE11).
\end{proof}

\begin{lemma}\label{te-real-gauss}
For any three lattices of the form (TE2), or any three
commensurable $\mathcal O$-stable self-dual trace lattices
with the same form, their triple scalar $\mu$ is real,
and hence $\mu^2=1$.
For an $E$-linear unitary map $g$ with $1-g$ invertible,
the Cayley Gauss sum in (TC31) is real.
\end{lemma}
\begin{proof}
For the first assertion choose $r>0$ carrying the sum
of the three lattices into their intersection.
Choose $m$ divisible by $2r^2$ and by the exponent of
the finite group $P/T$, and choose $\ell$ from (TE11).
Multiplication by $\ell$ permutes $P/T$: multiplication
by $-\bar\ell$ is its inverse, because $m(P/T)=0$.
Write $\ell\bar\ell+1=2r^2s$ with
$s\in\mathcal O_F$.  The scalar action has adjoint
$\bar\ell$, so (TE8) gives
\[
 \beta(\ell z,\ell y)+\beta(z,y)
                       =2\beta(rsz,ry)\in2\mathbb Z.
\]
Here both arguments on the right lie in the common
intersection, and $\beta$ is even integral there.
It follows that
$\psi(\ell z,\ell y)=\overline{\psi(z,y)}$.
The sum of $\psi$ is real, and so is $\mu$.

For the second assertion use the notation $A,C,\Gamma,q$
of (TC27), on an $\mathcal O$-stable lattice $\Lambda$.
The maps $A,C$ are $E$-linear.
Choose $r>0$ with $rC\Lambda\subset\Lambda$, and choose
$m$ divisible by $2r$ and the exponent of
$\Lambda/A\Gamma$.
Multiplication by the corresponding $\ell$ again
permutes this quotient.  Write
$\ell\bar\ell+1=2rs$, $s\in\mathcal O_F$.
Since $C^*=-C$, the exponent in
\[
 q(\ell x)q(x)=
 \exp\!\left(\pi i E_H(\tau(\ell\bar\ell+1)x,x)
                   -\pi H(C(\ell\bar\ell+1)x,x)\right)
 \tag{TE12}
\]
is $2\pi i$ times an integer.
Indeed the first $E_H$-pairing is even integral, and
$C(\ell\bar\ell+1)x=2s(rCx)\in2\Lambda$.
The second Hermitian value is purely imaginary, so its
imaginary part is even integral as well.
Thus $q(\ell x)=\overline{q(x)}$, proving reality of
its sum.
\end{proof}

\begin{proposition}\label{te-unitary-square}
Let $\Lambda$ be an $\mathcal O$-stable self-dual trace
lattice of even complex dimension $r$.  Let $g$ be
an $E$-linear unitary map.  The nonzero scalar defined by
\[
          F(g^{-1}z)=c(g,\Lambda)I_{\Lambda,g\Lambda}F(z)
                        \qquad(F\in\mathcal T(\Lambda,\alpha))
 \tag{TE13}
\]
satisfies
\[
                         c(g,\Lambda)^2=\det_{\mathbb C}g.
 \tag{TE14}
\]
\end{proposition}
\begin{proof}
The scalar exists by the one-dimensionality of the theta
spaces and the invariance of the norm semicharacter.
If $1-g$ is invertible, its nondegenerate Cayley Gauss
sum has absolute value $\sqrt N$ by (TC16), and is real
by Lemma~\ref{te-real-gauss}.  Its normalized square is
therefore one.  Equations (TC31) and (TC40) give
$c^2=(-1)^r\det g=\det g$.

For a general $g$, choose a scalar $z\in k$ of norm one
such that $z\ne1$ and $1-zg$ is invertible.
There are infinitely many choices before these finitely
many exclusions: with $s=\sqrt{-D}$, the values
$(1+ts)/(1-ts)$ for $t\in\mathbb Q$ are distinct and
have norm one.
Write $g=g_1g_2$, where $g_1=zg$ and $g_2=z^{-1}$.
Both factors have no fixed vectors.
Pullback commutes with normalized averaging after
transforming both lattices; this follows by reindexing
the finite sum and using preservation of $H$ and $\alpha$.
Consequently the scalar in (TE13) is the product of the
two factor scalars and a triple scalar for
$\Lambda,g_1\Lambda,g\Lambda$.
The latter has square one by Lemma~\ref{te-real-gauss}.
Squaring and using the first case proves (TE14).
The zero-dimensional case has scalar and determinant one.
\end{proof}

\subsection{The exact family and its three transformation laws}

Let $\theta$ be the fixed nonzero normalized theta
function for the rank-one trace lattice
$\lambda I\mathfrak a$, as constructed in
Sections~\ref{tc-section}--\ref{ts-section}.
Let $F_0$ be its product in $nn'$ independent entries,
and define
\[
 \begin{split}
 T(L,M)&=I_{\Lambda(L,M_0),\Lambda(L,M)}
                 I_{\Lambda(L_0,M_0),\Lambda(L,M_0)}F_0,\\
 b_\chi(L,M)&=\chi(\det(L_0,L))^{n'/2}
                             \chi(\det(M_0,M))^{-n/2},\\
 e_\chi(L,M)&=b_\chi(L,M)T(L,M).
 \end{split}
 \tag{TE15}
\]
Here $\det(L,L_1)$ is the relative determinant ideal.
The character $\chi$ is multiplicative on fractional
ideals $\mathfrak c$ with
$\mathfrak c\overline{\mathfrak c}=\mathcal O$ and obeys
\[
             \chi((u))=\prod_{\eta\in\Sigma}\eta(u)^{-1}
                          \qquad(u\bar u=1).
 \tag{TE16}
\]
These are exactly the character requirements used in
this assertion.  In particular $\chi$ never vanishes
on such an ideal.

\begin{theorem}\label{te-equivariance}
For positive even $n,n'$, the functions in (TE15)
satisfy all three identities
\[
 \begin{split}
 e_\chi(L,M_1)&=\chi(\det(M,M_1))^{-n/2}
                         I_{\Lambda(L,M),\Lambda(L,M_1)}e_\chi(L,M),\\
 e_\chi(L_1,M)&=\chi(\det(L,L_1))^{n'/2}
                         I_{\Lambda(L,M),\Lambda(L_1,M)}e_\chi(L,M),\\
 e_\chi(L,M)((\gamma'^{-1}z_\eta\gamma)_\eta)
                         &=e_\chi(\gamma L,\gamma'M)(z).
 \end{split}
 \tag{TE17}
\]
In the last equation $\gamma,\gamma'$ are unitary
matrices over $E$, and each matrix in the argument is
understood through its embedding $\eta$.
These identities also hold in dimension zero with the
empty-product conventions.
\end{theorem}
\begin{proof}
Proposition~\ref{te-exact-paths} gives the first two
identities for $T$ with the character factors removed.
The relative determinant is multiplicative in its
three lattice arguments: at any prime it is the
quotient of the two determinant lattices.  Its
Hermitian norm is one by self-duality.
Thus multiplication of the factors in (TE15) gives
the first two identities of (TE17).

For the last identity put
$D(\gamma)=\prod_{\eta\in\Sigma}\eta(\det_E\gamma)$.
First let only $\gamma$ act, at the standard pair.
Group the entries of $F_0$ into its $n'$ rows.
On each row the forward map of period lattices is
$x\mapsto x\gamma^{-1}$, whose complex determinant
is $D(\gamma)^{-1}$ and whose dimension $dn$ is even.
The operator in (TE13) and the normalized averaging
both factor over this orthogonal product: the finite
quotient is a product, and its positive square root
is the product of the positive square roots.
Proposition~\ref{te-unitary-square} therefore gives
\[
 F_0((z_\eta\gamma)_\eta)
                   =D(\gamma)^{-n'/2}T(\gamma L_0,M_0)(z).
 \tag{TE18}
\]
Likewise, grouping by columns gives
\[
 F_0((\gamma'^{-1}z_\eta)_\eta)
                    =D(\gamma')^{n/2}T(L_0,\gamma'M_0)(z).
 \tag{TE19}
\]
Here the forward map is $x\mapsto\gamma'x$, giving
the opposite determinant exponent.
Pullback commutes with averaging, as in the proof of
Proposition~\ref{te-unitary-square}.
The square and triangle identities in (TE10) therefore
extend (TE18)--(TE19) to every pair:
\[
 T(L,M)((\gamma'^{-1}z_\eta\gamma)_\eta)
       =D(\gamma)^{-n'/2}D(\gamma')^{n/2}
                                      T(\gamma L,\gamma'M)(z).
 \tag{TE20}
\]
To see explicitly why no path ambiguity remains,
transport the two defining edges in (TE15) by
$(\gamma,\gamma')$, precede them by the two edges
from $(L_0,M_0)$ to $(\gamma L_0,\gamma'M_0)$, and
interchange adjacent horizontal and vertical edges
using (TE10).  Consecutive edges in one direction
collapse by the corresponding triangle identity.
The result is exactly the two edges defining the
right side of (TE20).

Finally
\[
 \det(L_0,\gamma L)=(\det_E\gamma)\det(L_0,L).
\]
The same formula holds for $M$, and (TE16) gives
\[
 b_\chi(\gamma L,\gamma'M)
            =D(\gamma)^{-n'/2}D(\gamma')^{n/2}b_\chi(L,M).
\]
Substitution in (TE20) proves the last identity of
(TE17), with its exact scalar.
\end{proof}

This proves the covariance and the two moves required
for the actual normalized family.  It uses no choice
of signs for individual lattice moves.  For comparison,
the homogeneous quadratic function
$q(x)=e^{2\pi ix^2/3}$ on $\mathbb Z/3\mathbb Z$ has
sum $i\sqrt3$, so a general normalized quadratic Gauss
sum need not be one.  The local hyperbolic decomposition
and the norm congruence (TE11) are the specific arithmetic
facts that remove this obstruction here.  The measure
support at the selected primes and its use in incidence
maps remain separate assertions.


